\section{Integral Transforms and Representations}

The abstract solution $\ket{\psi(t)}$ acquires physical meaning only when
projected onto specific basis manifolds. We construct the explicit integral
transforms that map the angular momentum spectral amplitudes $\psi_{k\ell m}$ to
the observable representations of position, momentum, and energy.

\subsection{Position Space: The Spherical Bessel Transform}

The transition to coordinate space is defined by the transformation kernel
$\braket{\mathbf{r} | k, \ell, m}$. This kernel arises from the partial-wave
expansion of the plane wave, identifying the free spherical wave as the radial
eigenfunction.

\subsubsection{The Transformation Kernel}
We adopt the normalization convention consistent with the quantum Fourier
transform:
\begin{equation} \label{eq:pos_kernel}
	\braket{\mathbf{r} | k, \ell, m} = i^\ell \sqrt{\frac{2}{\pi}} j_\ell(kr)
	Y_{\ell m}(\Omega_r).
\end{equation}
The phase factor $i^\ell$ is essential for time-reversal consistency.

\subsubsection{Wavefunction Synthesis}
Projecting the spectral expansion Eq. (\ref{eq:general_solution_abstract}) onto
the position bra $\bra{\mathbf{r}}$ yields the spatial field:
\begin{align}
	\psi(\mathbf{r}, t) &= \sum_{\ell=0}^\infty \sum_{m=-\ell}^\ell i^\ell
	Y_{\ell m}(\Omega_r) \int_0^\infty \mathrm{d}k \, k^2 \,
	\underbrace{\sqrt{\frac{2}{\pi}} j_\ell(kr)}_{\text{Radial Basis}}
	\psi_{k\ell m}(t).
\end{align}
This expression identifies the radial wavefunction $R_{\ell}(r, t)$ as the
Spherical Bessel Transform (or Hankel transform of order $\ell + 1/2$) of the
spectral amplitude $\psi_{k\ell m}$. This is the spherical analog to the
standard Fourier transform, mapping the radial momentum domain $k$ to the radial
position domain $r$.

\subsection{Momentum Space: Harmonic Analysis on the Sphere}

In the linear momentum representation, the physics is simplified by the fact that
$\hat{\mathbf{p}}$ commutes with the free Hamiltonian. The angular momentum
states define a specific geometric partitioning of momentum space.

\subsubsection{Diagonal Representation}
The basis states $\ket{k, \ell, m}$ are eigenstates of the momentum magnitude
operator $\hat{p} = \sqrt{\hat{\mathbf{p}}^2}$. Consequently, the projection
kernel is sharply localized (distributional):
\begin{equation}
	\braket{\mathbf{p} | k, \ell, m} = \frac{\delta(p - \hbar k)}{p^2}
	Y_{\ell m}(\Omega_p).
\end{equation}
Here, $\Omega_p$ denotes the angles of the momentum vector.

\subsubsection{Momentum Distribution}
Substituting this kernel into the general solution collapses the radial integral
immediately:
\begin{align}
	\psi(\mathbf{p}, t) &= \int_0^\infty \mathrm{d}k \, k^2 \, \psi_{k\ell m}(t)
	\left[ \frac{\delta(p - \hbar k)}{p^2} Y_{\ell m}(\Omega_p) \right]
	\nonumber \\
											&= \frac{1}{\hbar^3} \psi_{p/\hbar, \ell, m}(t) Y_{\ell m}(\Omega_p).
\end{align}
This result confirms that the spectral coefficients $\psi_{k\ell m}$ are
physically equivalent to the momentum-space wavefunction projected onto
spherical harmonics. The time evolution is strictly a phase rotation, as
momentum states are stationary solutions of the free particle.

\subsection{Energy Representation and Density of States}

The energy representation is crucial for scattering theory (S-matrix formalism).
Since the spectrum of the free particle is continuous and degenerate, the
mapping involves the Density of States (DOS).

\subsubsection{Jacobian Transformation}
Using the dispersion relation $E = \frac{\hbar^2 k^2}{2m}$, we relate the
differential measures:
\begin{equation}
	\mathrm{d}k = \frac{\mathrm{d}k}{\mathrm{d}E} \mathrm{d}E = \frac{1}{\hbar}
	\sqrt{\frac{m}{2E}} \mathrm{d}E.
\end{equation}
The density of states factor is $\rho(E) \propto \sqrt{E}$.

\subsubsection{Spectral Evolution}
The wavefunction in the energy representation $\psi(E, \ell, m)$ is related to
the $k$-space amplitude by preserving the probability norm $\int |\psi|^2
\mathrm{d}k = \int |\phi|^2 \mathrm{d}E$:
\begin{equation}
	\psi(E, \ell, m; t) = \left( \frac{m k(E)}{\hbar^2} \right)^{1/2}
	\psi_{k(E), \ell, m}(0) e^{-iEt/\hbar}.
\end{equation}
This representation trivializes the dynamics to a linear phase evolution,
explicitly demonstrating that the angular momentum modes are the stationary
states of the system.
